\documentclass[11pt]{article}
\usepackage[margin=1in]{geometry}
\usepackage{amsmath,amssymb,booktabs,hyperref,longtable,xcolor,graphicx}
\usepackage[T1]{fontenc}
\title{Independent replication of\\
       ``Implementation of the Quantum Fourier Transform''\\
       (Weinstein, Lloyd, Cory, quant-ph/9906059)\\[0.3em]
       \large QC-200 wave, replication host = CherryRd}
\author{Ollie (agent), for R. Stevens}
\date{2026-07-05}

\begin{document}
\maketitle

\begin{abstract}
We independently replicate the algorithmic content of Weinstein, Lloyd, and
Cory's 1999 paper on implementing the Quantum Fourier Transform (QFT) on a
3-qubit NMR quantum computer. The paper's testable core is the
Coppersmith decomposition of the QFT into $n$ Hadamards and $n(n-1)/2$
controlled-phase gates, and the explicit $\text{QFT}_4$ matrix given in
its Eq.~4. Using Qiskit~2.5.0 we build the Coppersmith circuit for
$n=3,4,5$ and verify (i)~the gate count is exact
($H=n$, $\text{CP}=n(n-1)/2$), (ii)~the output state matches the analytic
QFT amplitudes on every computational-basis input to
$\le 5.6\times10^{-15}$, (iii)~the $n=2$ operator matches paper
Eq.~4 to $1.15\times10^{-16}$, and (iv)~our circuit agrees with
\texttt{qiskit.circuit.library.QFT} to zero after removing a global
phase. The paper's NMR hardware fidelity numbers ($F=87\%$ overall,
$>98\%$ per gate) are hardware claims and are not reproducible on an
ideal statevector simulator; we note the correct interpretation and
scope. \textbf{Verdict: REPLICATED} for the paper's algorithmic and
computational claims; the hardware claims are marked
\emph{OUT-OF-SCOPE-BY-DESIGN}.
\end{abstract}

\section{Paper summary}

Weinstein, Lloyd, and Cory (MIT, 1999; arXiv:quant-ph/9906059v1)
implement the QFT on the three C-13 spins of an alanine molecule at
9.4~T on a liquid-state NMR spectrometer. The QFT is defined
(paper Eq.~1) as
\[
\text{QFT}_q\,|x\rangle \to \frac{1}{\sqrt{q}}\sum_{p=0}^{q-1}
   e^{2\pi i a p / q}\,|p\rangle,
\]
and for the two-qubit ($q=4$) case reduces to the explicit unitary in
paper Eq.~4. The paper adopts the Coppersmith decomposition of the QFT
(paper Eqs.~5--7): $A_j$ is a Hadamard on qubit $j$, and $B_{jk}$ is a
controlled-phase gate with angle $\theta_{jk} = \pi/2^{k-j}$ on the
$(j,k)$ pair. The composite gate $B_{j,j+1} B_{j,j+2}\ldots B_{j,L-1}
A_j$ is applied on the ``lead bit'' $j = L-1$, then the sequence is
repeated with $j$ decremented to $0$; this gives an $O(n^2)$-two-qubit-gate
QFT. The paper then translates each Coppersmith gate into an NMR pulse
sequence for a $J$-coupled 3-spin system (paper Eqs.~9--11), and
reports:
\begin{itemize}
  \item overall NMR fidelity $F = 87\%$ using the Uhlmann-like measure
        in paper Eq.~12;
  \item after renormalizing $\rho_\text{exp}$ to compensate for
        decoherence attenuation, per-gate fidelity $>98\%$ over the six
        gates of the 3-qubit QFT.
\end{itemize}

\section{Claims table}

\begin{longtable}{p{0.5cm} p{5.5cm} p{2.2cm} p{2cm} p{4cm}}
\toprule
\# & Claim & Type & Testable in ideal sim? & Tested here? \\
\midrule
\endhead
C1 & QFT decomposes into $A_j = H$ (paper Eq.~5) and
     $B_{jk} = \text{CP}(\pi/2^{k-j})$ (paper Eq.~6). &
     algorithmic & YES & YES \\
C2 & For an $L$-qubit register the Coppersmith QFT uses
     $n$ Hadamards and $n(n-1)/2$ controlled-phase gates
     ($O(n^2)$ total two-qubit gates). &
     complexity & YES & YES (n=3,4,5) \\
C3 & The QFT satisfies paper Eq.~1 for all computational-basis
     inputs. & correctness & YES & YES (n=3,4,5) \\
C4 & $\text{QFT}_4$ equals the explicit matrix in paper Eq.~4. &
     correctness & YES & YES \\
C5 & Alanine C-13 3-spin NMR realization achieves overall
     $F = 87\%$. & hardware & NO & OUT-OF-SCOPE (needs a real
     spectrometer) \\
C6 & After renormalizing $\rho_\text{exp}$ to compensate for
     decoherence, per-gate NMR fidelity is $>98\%$ over 6 gates. &
     hardware & NO & OUT-OF-SCOPE \\
\bottomrule
\end{longtable}

\section{Method}

\subsection{Environment}
\begin{itemize}
  \item Host: \texttt{CherryRd} (macOS; x86\_64), Python 3.14.6.
  \item Qiskit 2.5.0, NumPy 2.5.1, MarkItDown for the marker-side
        extraction (see failure analysis for why Marker/Nougat
        proper were not used).
  \item All simulation is exact statevector; no NMR / density-matrix
        / open-system simulation.
\end{itemize}

\subsection{Circuit construction (Coppersmith)}
For $j = n-1,\, n-2,\, \ldots,\, 0$ apply
\[
  H_j,\; \text{CP}(\pi/2^{j-k})_{k \to j}\quad\text{for } k=j-1,\ldots,0.
\]
Optionally append the bit-reversal SWAP. This is a direct transcription
of paper Eq.~7 into Qiskit's little-endian convention. See
\texttt{report/evidence/qft\_replication.py}, function
\texttt{coppersmith\_qft}.

\subsection{Verification procedures}
\begin{enumerate}
  \item \textbf{Gate count} (\texttt{verify\_gate\_count\_claim}):
        build the circuit without the trailing SWAP, count \texttt{h}
        and \texttt{cp} operations, compare to $n$ and $n(n-1)/2$.
  \item \textbf{Analytic correctness} (\texttt{verify\_correctness}):
        for every $x\in\{0,\ldots,2^n-1\}$ prepare $|x\rangle$, apply
        the circuit, and compare the resulting statevector to
        $\frac{1}{\sqrt{2^n}}\sum_y e^{2\pi i x y / 2^n}|y\rangle$.
  \item \textbf{Explicit matrix} (\texttt{verify\_paper\_eq4\_matrix}):
        build the full 4x4 operator for $n=2$ and compare to paper Eq.~4
        entry-by-entry.
  \item \textbf{Sanity against Qiskit's own QFT}
        (\texttt{compare\_to\_qiskit\_builtin}): fold out a possible
        global phase between our operator and
        \texttt{qiskit.circuit.library.QFT}.
  \item \textbf{Anchor cases} (\texttt{anchor\_test\_case}): print
        amplitudes for $x=0,1,3,7$ at $n=3$.
\end{enumerate}
Reproduce with:
\begin{verbatim}
cd QC-quant-ph-9906059-quantum-fourier-transform-implementation
python3 -m venv .venv && source .venv/bin/activate
pip install qiskit numpy markitdown
python report/evidence/qft_replication.py
\end{verbatim}

\section{Results vs. paper}

\subsection{Gate count (C2)}
\begin{center}
\begin{tabular}{cccccc}
\toprule
$n$ & $H$ measured & $H$ expected & CP measured & CP expected & match? \\
\midrule
3 & 3 & 3 & 3 & 3 & YES \\
4 & 4 & 4 & 6 & 6 & YES \\
5 & 5 & 5 & 10 & 10 & YES \\
\bottomrule
\end{tabular}
\end{center}
Formula $n(n-1)/2$ for controlled-phase count is exactly satisfied.

\subsection{Analytic correctness (C3)}
Max $\ell_\infty$ error over all $2^n$ computational-basis inputs:
\begin{center}
\begin{tabular}{cccc}
\toprule
$n$ & inputs checked & max amplitude error vs analytic & to machine precision? \\
\midrule
3 & 8  & $1.42\times10^{-15}$ & YES \\
4 & 16 & $3.78\times10^{-15}$ & YES \\
5 & 32 & $5.62\times10^{-15}$ & YES \\
\bottomrule
\end{tabular}
\end{center}

\subsection{Paper Eq.~4 matrix (C4)}
Our $\text{QFT}_4$ operator matches the explicit matrix in paper Eq.~4
to $1.15\times10^{-16}$ in every entry.

\subsection{Sanity vs Qiskit built-in QFT}
For each $n\in\{3,4,5\}$, our Coppersmith operator equals
\texttt{qiskit.circuit.library.QFT} to floating-point zero after
absorbing a global phase; phase variation across nonzero entries is
zero to floating-point.

\subsection{3-qubit anchor cases}
For $x=0,1,3,7$ at $n=3$: max amplitude error vs
$\frac{1}{\sqrt8}\sum_y e^{2\pi i xy/8}|y\rangle$ is
$\le 5\times10^{-16}$ in every case (see
\texttt{report/evidence/results.json},
key \texttt{anchor\_3qubit\_selected\_inputs}).

\subsection{Hardware claims (C5, C6)}
Not tested. These are NMR spectrometer measurements on an alanine
sample; they require an actual 9.4~T NMR magnet and a $J$-coupled
$^{13}$C sample. We note that the paper's per-gate 98\% comes from a
renormalization argument (isotropic attenuation of the deviation
density matrix) that is not obviously exact when $J_{13}$ is 30--45x
smaller than $J_{12}/J_{23}$, so refocusing intervals approach $T_2$;
we flag this in Open Question Q2.

\section{Verdict}

\textbf{REPLICATED} for the paper's algorithmic and computational
content:
\begin{itemize}
  \item C1 (Coppersmith decomposition): reproduced exactly.
  \item C2 (gate count $n(n-1)/2$ CP + $n$ H): reproduced exactly
        for $n=3,4,5$.
  \item C3 (analytic correctness on all basis inputs): reproduced to
        machine precision.
  \item C4 (paper's Eq.~4 QFT$_4$ matrix): reproduced to
        $10^{-16}$.
\end{itemize}
Hardware claims C5 and C6 are \emph{OUT-OF-SCOPE-BY-DESIGN} for a
statevector replication and are not evaluated as MATCH/MISMATCH.

\section{Open Questions}

\begin{enumerate}
  \item[Q1.] How well does the Coppersmith QFT circuit, dropped into
        the paper's Eq.~11 pulse program, reproduce the reported
        87\% overall NMR fidelity when composed with a realistic
        amplitude-damping + phase-damping channel calibrated to
        $T_1 = 1.56\,$s and $T_2 = 420\,$ms on alanine C-13 at
        100.6~MHz, rather than an isotropic depolarizing channel?
        (Motivated by: our ideal statevector rules out the algorithmic
        decomposition as a source of the 13-point fidelity gap, but we
        did not run an open-system simulation.)
  \item[Q2.] The paper's ``renormalize $\rho_\text{exp}$ to lift $F$
        from 87\% to 98\%/gate'' step assumes decoherence acts as an
        isotropic scalar attenuation on the deviation density matrix.
        Under which noise channels is that actually valid, and does it
        hold when $J_{13} = 1.2\,$Hz forces $1/(8 J_{13}) \sim
        100\,$ms refocusing intervals comparable to $T_2$?
  \item[Q3.] The paper avoids an explicit SWAP by ``reordering the bits
        at the appropriate interval''. The exact mapping from
        (carbonyl, $C_\alpha$, $C_\beta$) to Qiskit qubit indices is
        not printed, blocking a bit-for-bit reproduction of Fig.~1's
        A/B/C spectrum-to-qubit assignment.
  \item[Q4.] Up to what $n$ is the Coppersmith $n(n-1)/2$ controlled-
        phase count the actual optimum for the QFT on an all-to-all
        register, before approximate-QFT truncation of small-angle
        $B$-gates or restricted connectivity change the optimum?
  \item[Q5.] Given our exact ideal-statevector result at $n=3,4,5$, on
        modern superconducting/trapped-ion hardware with two-qubit gate
        error $\sim 10^{-3}$, what is the crossover $n^\star$ at which
        the full Coppersmith QFT becomes indistinguishable, via
        Eq.~12's $F$, from a Haar-random unitary?
\end{enumerate}

See \texttt{report/open\_questions.json} for the machine-readable
version with \texttt{next\_steps}.

\section{Artifacts}
See \texttt{report/artifacts\_summary.md} for the full 8-artifact
inventory and \texttt{report/failure\_analysis.md} for the residual
gaps (Marker/Nougat model weights not on this host; hardware claims
by-design out of scope).

\end{document}
